\subsection{WEIGHTS, RADICAL WEIGHTS}

Let \(l = \dim V\) . Denote by \(Q(R)\) the subgroup of \(V\) generated by \(R\) ; the elements of \(Q(R)\) are called the radical weights of \(R\) . By Th. 3 of no. 6, \(Q(R)\) is a discrete subgroup of \(V\) of rank \(l\) , and every basis of \(R\) is a basis of \(Q(R)\) .

Similarly, the group \(\mathrm{Q}(\mathsf{R}^{\cdot})\) is a discrete subgroup of \(V^{*}\) of rank l.

PROPOSITION 26. The set of \(x \in V\) such that \(\langle x, y^* \rangle \in \mathbf{Z}\) for all \(y^* \in Q(R^*)\) (or, equivalently, for all \(y^* \in R^*)\) is a discrete subgroup \(G\) of \(V\) containing \(Q(R)\) . If \(B'\) is a basis of \(R^*\) , the basis of \(V\) dual to \(B'\) is a basis of \(G\) .

Let \(x \in V\) . The following three properties are equivalent:

\begin{enumerate}
\def\labelenumi{(\roman{enumi})}
\item
  \(\langle x, y^{*} \rangle \in \mathbf{Z}\) for all \(y^{*} \in \mathrm{Q}(\mathrm{R}^{\vee})\) ;
\item
  \(\langle x, y^{*} \rangle \in Z\) for all \(y^{*} \in B'\) ;
\item
  the coordinates of \(x\) with respect to the dual basis of \(\mathbf{B}'\) are in \(\mathbf{Z}\) .
\end{enumerate}

We deduce from this that the basis dual to \(B'\) is a basis of \(G\) . On the other hand, \((RS_{\mathrm{III}})\) proves that \(R \subseteq G\) , so \(Q(R) \subseteq G\) .

The group G of Prop. 26 is denoted by P(R), and its elements are called the weights of R. We can also consider the group P(R') of weights of R'.

By Algebra, Chap. VII, 2nd edn., § 4, no. 8,

\[
\mathrm{P} (\mathrm{R}) / \mathrm{Q} (\mathrm{R}), \quad \mathrm{P} (\mathrm{R} ^ {\cdot}) / \mathrm{Q} (\mathrm{R} ^ {\cdot})
\]

are finite groups in duality over \(\mathbf{Q} / \mathbf{Z}\) , and hence are isomorphic. The common order of these two groups is called the connection index of R (or of R \(^{\sim}\) ).

If R is a direct sum of root systems \(R_{i}\) , the group \(Q(R)\) (resp. \(P(R)\) ) is identified canonically with the direct sum of the \(Q(R_{i})\) (resp. \(P(R_{i})\) ).

PROPOSITION 27. Let \(\mathbb{R}_1\) be a subset of \(\mathbb{R}\) , \(\mathbb{Q}_1\) the subgroup of \(\mathbb{Q}(\mathbb{R})\) generated by \(\mathbb{R}_1\) , and \(\mathbb{W}_1\) the subgroup of \(\mathbb{W}(\mathbb{R})\) generated by the \(s_\alpha\) ( \(\alpha \in \mathbb{R}_1\) ). If \(p \in \mathrm{P}(\mathbb{R})\) and \(w \in \mathbb{W}_1\) , then \(p - w(p) \in \mathbb{Q}_1\) .

If \(w = s_{\alpha}\) with \(\alpha \in \mathbb{R}_1\) , then

\[
p - w (p) = \langle p, \alpha^ {\prime} \rangle \alpha \in \mathbf {Z} \alpha \subseteq \mathrm{Q} _ {1}.
\]

If \(w = s_{\alpha_1}s_{\alpha_2}\ldots s_{\alpha_r}\) , with \(\alpha_{1},\ldots ,\alpha_{r}\in \mathbb{R}_{1}\) , it is still true that \(p - w(p)\in \mathbf{Q}_1\) , as we see by induction on \(r\) .

The group A(R) leaves P(R) and Q(R) invariant, and hence acts on the quotient P(R)/Q(R). By Prop. 27, the group W(R) acts trivially on P(R)/Q(R). Passing to the quotient, we see that the quotient group A(R)/W(R) (cf.~Prop. 16, no. 5) acts canonically on P(R)/Q(R).

